
\chapter{Teamorganisation und Projekterfahrungen}

\section{Teamorganisation}
Unsere Zusammenarbeit im Team war eine Mischung aus klarer Aufgabenverteilung und engem Austausch. Für schnelle Rückfragen, spontane Abstimmungen und das Klären von Problemen haben wir vor allem Discord genutzt. Dort haben wir kurze Statusupdates ausgetauscht, offene Punkte gesammelt und bei Bedarf direkt in Calls besprochen. Die eigentliche Projektdokumentation haben wir parallel in Overleaf geschrieben, sodass Abschnitte direkt gemeinsam bearbeitet, kommentiert und zusammengeführt werden konnten. Google Sheets diente uns als Aufgabenübersicht mit Zuständigkeiten, offenen Punkten und Zwischenständen, während wir in Google Docs eher lose Notizen für die 15 Services gesammelt haben. Für Code, Konfigurationsstände und nachvollziehbare Änderungen haben wir GitHub genutzt.

Trotz der fachlichen Aufteilung haben wir das Projekt nicht als drei voneinander getrennte Teilaufgaben bearbeitet. Gemeinsame Basisdienste wie PostgreSQL, \ac{NATS}, WireGuard sowie Logging und Monitoring wurden bewusst als geteilte Infrastruktur verstanden. Das hat Doppelarbeit vermieden, bedeutete aber auch, dass wir Schnittstellen sauber abstimmen mussten. Gerade an diesen Übergängen hat sich gezeigt, wie wichtig regelmäßige Kommunikation für ein konsistentes Gesamtsystem war.

\section{Aufgabenverteilung}
Wir haben die Arbeit entlang klarer Schwerpunkte verteilt. Luca übernahm die Infrastruktur- und Self-Hosted-Themen. Dazu gehörten das Serverhosting, der Aufbau des Self-Hosted-Setups und viele der grundlegenden Architekturentscheidungen auf Plattformebene.

Florian verantwortete den Bitcoin-Bereich. Sein Schwerpunkt lag auf der Konzeption der Hot- und Cold-Wallet-Architektur, dem Transaktionssystem sowie den sicherheitsrelevanten Abläufen rund um Signierung, Schlüsselverwaltung und Automatisierung.

Marie übernahm die KI-Evaluation und große Teile der Dokumentationskoordination. Dazu gehörten die Bewertung von KI-Unterstützung, die Pflege der Dokumentation und die Ausarbeitung allgemeiner Kapitel.

Wichtig war für uns dabei, dass die Zuständigkeiten zwar klar, aber nicht starr waren. Wenn aus einer Infrastrukturentscheidung Folgen für die Bitcoin-Komponente entstanden oder Dokumentation und Implementierung auseinanderliefen, haben wir das gemeinsam besprochen. 


\section{Herausforderungen}
\subsection{Self-Hosted-Setup}
Eine der ersten größeren Herausforderungen für uns im Infrastrukturteil war die Frage, welche Architektur für das Projekt wirklich sinnvoll ist. Anfangs war Kubernetes geplant, im Verlauf der Arbeit hat sich Docker Compose für unseren Rahmen aber als die stabilere und besser beherrschbare Lösung herausgestellt.

Dazu kamen Plattformthemen rund um NixOS und Isolationsmechanismen. Ein NixOS-Upgrade kurz vor der Abgabe erzeugte zusätzlichen Anpassungsdruck, während Technologien wie Kata-Container oder \textit{gVisor} nicht in allen Konstellationen so stabil liefen, wie wir es uns erhofft hatten. Zusätzlich erschwerten auch die Rahmenbedingungen des Projekts die Arbeit: Durch die parallele Belastung aus Beruf und Studium sowie die unterschiedlichen Erfahrungsstände im Team dauerten Abstimmung, Tests und belastbare Entscheidungen oft länger als geplant.

\subsection{Bitcoin-Setup}
Im Bitcoin-Teil haben wir besonders gemerkt, wie schwierig sicherheitskritische Eigenentwicklungen ohne belastbare Referenzen sind. Mehrere Architekturansätze haben wir zunächst aufgebaut und später wieder verworfen, weil sie zwar interessant wirkten, den Gesamtaufbau aber unnötig verkomplizierten. Zusätzlich hat uns die lückenhafte Dokumentation der Bibliothek \texttt{embit} gebremst, vor allem bei der Wallet-Verwaltung und beim automatisierten Signieren.

Aufwendig war auch der Versuch, dem Mandanten möglichst wenig manuelle Arbeit zu überlassen. Dadurch entstand viel zusätzliche Logik in Skripten und Abläufen. Gleichzeitig setzte uns die Laborumgebung Grenzen: Der Austausch über QR-Codes ließ sich in der Proxmox-basierten Umgebung nicht realistisch nachstellen, sodass wir auf andere Wege ausweichen mussten. Hinzu kamen Probleme mit Sparrow unter dem stabilen NixOS-Zweig sowie der generell hohe Rebuild- und Testaufwand, den der deklarative NixOS-Ansatz mit sich bringt.

Besonders deutlich wurde für uns der operative Teil von Sicherheit beim Zugriff auf das TPM aus einem unprivilegierten Container. Hier ging es nicht nur darum, dass etwas grundsätzlich funktioniert, sondern dass minimale Rechte, reproduzierbare Konfiguration und realer Hardwarezugriff gleichzeitig zusammenpassen. Gerade an solchen Stellen wurde sichtbar, dass Sicherheit nicht nur aus Kryptographie besteht, sondern auch aus sauberen Betriebsabläufen.

\subsection{KI-Evaluation}
Im Bereich \acs{KI} war eine der ersten Herausforderungen die große Auswahl an verfügbaren Modellen und Werkzeugen. Gerade am Anfang war es nicht leicht einzuschätzen, welches Modell für welche Aufgabe am besten geeignet ist. Deshalb wurde pragmatisch entschieden, zunächst mit den Modellen zu arbeiten, die im Alltag ohnehin schon genutzt wurden und deren Verhalten bereits etwas bekannt war.

Eine weitere Herausforderung war, dass die Idee, Codex gezielt in die Evaluation einzubeziehen, erst relativ spät entstand. Dadurch mussten zusätzliche Testläufe geplant und eingerichtet werden. Das war deutlich aufwendiger als eine reine Chat-Auswertung, weil Codex nicht nur Antworten schreiben, sondern tatsächlich in einer lokalen Testumgebung arbeiten sollte und diverse Dateien erstellt hat, die ebenfalls analysiert werden mussten. Die Läufe dauerten entsprechend lange und konnten nicht einfach nebenbei durchgeführt werden.

Schwierig war außerdem die Umsetzung von Szenario B. Dort sollte ein fachkundiger Nutzer simuliert werden. Dafür wurde Wissen aus den beiden anderen Projektteilen benötigt. Das ließ sich nur eingeschränkt vollständig nachbilden, weil die Person, die den \acs{KI}-Teil bearbeitet hat, nicht in gleicher Tiefe an beiden anderen Teilen beteiligt war. Gleichzeitig ließ sich die Arbeit auch nicht einfach aufteilen, weil Codex auf einer lokalen \acs{VM} gearbeitet hat und viele Entscheidungen direkt während des laufenden Tests getroffen werden mussten.

Auch die spätere Auswertung war dadurch nicht ganz trivial und musste durch die beiden anderen Gruppenmitglieder stark unterstützt werden. Für den Bitcoin-Teil brauchte man Detailwissen zu Wallets, PSBT, Watch-only-Konzepten, Regtest und Schlüsseltrennung. Für die Infrastruktur waren dagegen Kenntnisse in den Bereichen Netzwerktrennung, Tor-Zugriffen, Backups, Monitoring und Härtung wichtig. Entweder hätten sich die anderen Gruppenmitglieder zusätzlich in den \acs{KI}-Teil einarbeiten müssen, oder die Person im \acs{KI}-Teil hätte sich sehr tief in beide Fachbereiche einarbeiten müssen. Beides wäre wegen des Umfangs und des Zeitaufwands nur schwer realistisch gewesen.

Das erklärt auch, warum Szenario B nicht in allen Bereichen ideal vorbereitet war. Der Anspruch war zwar, einen Expertennutzer nachzuahmen, aber dieses Expertenwissen musste aus mehreren Projektteilen zusammengeführt werden. Gerade im Bitcoin-Bereich hätten einige Anforderungen von Anfang an noch genauer formuliert werden können.


\section{Learnings}
Ein wichtiges Learning für uns war, dass nicht jede technisch mögliche Erweiterung am Ende einen echten Mehrwert bringt. Mehrere Ansätze waren auf dem Papier interessant, hätten das System in der Praxis aber unnötig komplexer und angreifbarer gemacht.

Für den Infrastrukturteil haben wir vor allem gelernt, wie wichtig frühe und klare Entscheidungskriterien sind. Sie helfen dabei, Tool-Bias zu vermeiden und aus vielen Optionen schneller zu einer tragfähigen Lösung zu kommen. Im Rückblick hätten uns einige frühere Festlegungen Zeit und Umwege erspart.

Im Bitcoin-Bereich war die wichtigste Erkenntnis für uns, dass der Eigenbau einer vollständigen Custody-Lösung erheblich komplexer ist, als es am Anfang wirkt. Bewährte Komponenten wie Bitcoin Core, Sparrow und Hardware-Wallets reduzieren nicht nur den Implementierungsaufwand, sondern können die Sicherheit auch erhöhen, weil sie etablierte und nachvollziehbare Funktionen mitbringen. Für uns war das ein klarer Hinweis darauf, dass Sicherheit nicht automatisch durch möglichst viel Eigenentwicklung entsteht.

Auch organisatorisch haben wir mitgenommen, dass klare Zuständigkeiten und gemeinsame Werkzeuge gut zusammenpassen. Gerade weil technische Umsetzung und Dokumentation parallel liefen, hat uns diese Kombination geholfen, als Team konsistent zu arbeiten.


\section{Feedback und Anmerkung zur Gewichtung}

Aus unserer Sicht war der Gesamtumfang der Aufgabenstellung deutlich zu hoch. Besonders in Kombination mit Aufgabenteil 1, der für sich genommen bereits auf mehrere Wochen ausgelegt war, entstand insgesamt eine sehr große Arbeitsbelastung. Gegen Ende des Moduls war bei allen drei Gruppenmitgliedern spürbar die Luft raus. Die Menge an Aufgaben, die technische Tiefe und der Zeitdruck führten dazu, dass wir teilweise überfordert waren und das Modul insgesamt stark belastend wurde.

Beide Aufgabenteile waren inhaltlich sehr interessant. Gerade deshalb wäre es aus unserer Sicht sinnvoll, den Umfang im nächsten Semester stärker zu fokussieren. Unser Vorschlag wäre, Aufgabenteil 1 zu streichen und stattdessen das gesamte Semester für den zweiten, praktischen Teil zu nutzen. Dadurch bliebe mehr Zeit, sich gründlich in die Themen einzuarbeiten, Dinge praktisch auszuprobieren und die Ergebnisse sauberer zu dokumentieren.

Zusätzlich wäre es aus unserer Sicht hilfreich, den zweiten Aufgabenteil selbst etwas zu reduzieren. Statt drei großer Teilbereiche könnte ein Bereich gestrichen oder der erwartete Umfang deutlich begrenzt werden. Dann hätten alle Gruppenmitglieder mehr Zeit gehabt, sich nicht nur mit dem eigenen Abschnitt, sondern auch mit den anderen Themen zu beschäftigen. So hätte jede Person aus mehreren Bereichen etwas lernen und praktische Erfahrungen sammeln können.

Durch den hohen Zeitdruck mussten wir die Arbeit am Ende stark aufteilen: Jede Person hat vor allem ihren eigenen Bereich fertiggestellt, während die anderen dort kaum praktisch mitarbeiten konnten. Das war schade, weil alle Bereiche fachlich spannend waren. Mit weniger Umfang wäre aus unserer Sicht nicht nur die Belastung geringer gewesen, sondern auch der Lerneffekt für alle Gruppenmitglieder deutlich größer.

Vor diesem Hintergrund halten wir eine zu starke Gewichtung einzelner Teilbereiche innerhalb des Projekts nur für bedingt sinnvoll. Grundsätzlich würden wir die drei Hauptbereiche gleich gewichten, da alle Teile für das Gesamtsystem notwendig waren und stark voneinander abhingen. Infrastruktur, Bitcoin-Komponente und \acs{KI}-Evaluation erfüllten jeweils unterschiedliche Funktionen und lassen sich deshalb nur schwer direkt miteinander vergleichen.

Wenn überhaupt eine unterschiedliche Gewichtung vorgenommen wird, müsste aus unserer Sicht eher der Bitcoin-Teil höher gewichtet werden. Dort war der Anteil an eigener Konzeption, Erklärung und sicherheitskritischer Abwägung besonders hoch. Anders als bei vielen Infrastrukturthemen existieren für den konkreten Aufbau einer solchen mandantenbezogenen Bitcoin- und Transaktionsarchitektur nur begrenzt direkt übertragbare Referenzen. Viele Entscheidungen mussten daher selbst hergeleitet, begründet und auf ihre praktischen sowie sicherheitstechnischen Folgen geprüft werden.

Gerade dieser Teil erforderte besonders viel eigenständige Arbeit, auch wenn sich das nicht immer unmittelbar in sichtbaren Konfigurationsdateien oder im Seitenumfang widerspiegelt. Deshalb sollte bei der Bewertung nicht nur betrachtet werden, wie viel technisch umgesetzt oder dokumentiert wurde, sondern auch, wie viel konzeptionelle und sicherheitskritische Eigenleistung in den jeweiligen Teilbereichen steckt.
